% source: https://tex.stackexchange.com/a/765462 (question 198211, score 5, block 1)
% Source - https://tex.stackexchange.com/a/765452
% Posted by pascal974
% Retrieved 2026-08-15, License - CC BY-SA 4.0

 % !TeX program = lualatex
 \documentclass{standalone}
 \usepackage[3d]{luadraw}
 %
 \begin{document}
 \begin{luadraw}{name=truncated_cone}
  local ld = luadraw 
 local pt3d = ld.pt3d
 local M = pt3d.M
 local Origin, vecI, vecJ, vecK = pt3d.Origin, pt3d.vecI, pt3d.vecJ, pt3d.vecK
 --
 local g = ld.graph3d:new{
    window3d = {-6,6,-4,4,-1,6},
    viewdir={"yz",0.3,60},
    window={-2,7,-5,5}, 
    size={10,10}
 }
 ld.Hiddenlines = true
 ld.Hiddenlinestyle = "dashed"
 --
 g:Setmatrix3d({Origin, vecI, vecK, vecJ}) -- y-axis <-> z_axis
 --
 local R, r, h, d = 2, 1, 3, 1.5
 local A = M(0,0,d+h)-- ⟨A⟩ est un point 3D représentant le centre de la face de rayon ⟨R⟩
 local B = M(0,0,d)

g:Dscene3d(
    g:addPoly( ld.frustum(A, R, r, vecK, B, 100), {color="green"}),
    g:addAxes(M(0,0,0), {arrows=1, width=8})
    )
 g:Dfrustum(A, R, r, vecK, B)
  g:Show()
 \end{luadraw}
 \end{document} 
